\section{AI 驱动的科学与药物发现}

\subsection{AGI 对科学的革命性影响}

在所有 AGI 可能改变的领域中，Hassabis 最为激动的是\textbf{科学发现和医学}。他认为 AGI 有潜力彻底改变科学研究的范式——从人类主导的假说驱动型研究，转向 AI 辅助的大规模模拟与发现。

\begin{importantbox}{Hassabis 的终极目标}
Hassabis 的核心使命不仅仅是构建 AGI，而是\textbf{用 AGI 来推进科学、治愈疾病}。他明确表示，自己的遗产（legacy）目标是：推进人类科学认知、攻克重大疾病。

这也是他同时领导 Google DeepMind 和 Isomorphic Labs 的原因——前者追求通用智能，后者将 AI 直接应用于药物发现。
\end{importantbox}

\subsection{Isomorphic Labs：AI 药物发现}

Hassabis 创立了 Isomorphic Labs，专注于用 AI 革新药物发现流程。目标是建立一个\textbf{通用的药物设计平台}（general-purpose drug design platform），最终目标是\textbf{攻克癌症}。

药物发现的 AI 化路径包括：
\begin{enumerate}
    \item \textbf{化合物设计}：AI 生成候选药物分子
    \item \textbf{毒性预防}：在设计阶段就预测和规避毒性问题
    \item \textbf{代谢模拟}：用 AI 模拟药物在人体内的代谢过程
    \item \textbf{回溯验证}：将模拟结果与已知数据进行交叉验证
\end{enumerate}

预计时间线：5--10 年内实现变革性的成果。

\begin{importantbox}{跳过动物实验：药物审批加速的两步策略}
Hassabis 提出了一种大胆的药物审批加速方案：

\textbf{Step 1}：AI 设计药物 $\to$ 模拟人体代谢过程 $\to$ 与已知数据回溯验证

\textbf{Step 2}：如果模拟精度足够高，可以\textbf{跳过部分动物实验}，直接进入人体临床试验

这需要监管机构的适应和配合——如果 AI 的预测准确率远超动物模型，继续坚持传统的动物测试流程就变成了不必要的时间浪费。
\end{importantbox}

\subsection{AlphaFold 的启示}

AlphaFold 是 Hassabis 关于``AI 驱动科学发现''愿景的最佳范例。该系统解决了生物学界 50 年来的核心挑战——蛋白质结构预测问题，为此 Hassabis 获得了 2024 年诺贝尔化学奖。

\begin{knowledgebox}{从 AlphaFold 到通用药物设计}
AlphaFold 的成功证明了一个关键命题：\textbf{AI 可以在科学问题上达到甚至超越人类专家的水平}。但蛋白质折叠只是药物发现链条中的一环。从``预测蛋白质结构''到``设计有效且安全的药物''还需要跨越巨大的鸿沟：理解蛋白质-配体相互作用、预测药代动力学、评估脱靶效应等。

Isomorphic Labs 的使命正是从 AlphaFold 的成功出发，向着完整的药物设计流水线延伸。
\end{knowledgebox}

\subsection{监管必须适应}

Hassabis 强调，\textbf{监管流程必须适应 AI 时代的速度}。如果 AI 能在几天内完成传统上需要数年的药物设计和安全性评估，但审批流程仍然需要 10 年，那么技术突破的价值就会被制度瓶颈严重稀释。

\begin{warningbox}{制度瓶颈可能比技术瓶颈更致命}
最大的风险不是 AI 做不到，而是\textbf{即使 AI 做到了，制度系统也跟不上}。药物审批、临床试验规范、国际协调机制——这些人类制度的进化速度远慢于 AI 技术的进步速度。

如果我们不主动改革监管流程，AI 发现的救命药物可能会在审批流水线上等待数年，而在此期间患者在继续死去。
\end{warningbox}

\subsection{本章小结}

AI 驱动的药物发现是 Hassabis 最核心的使命。Isomorphic Labs 致力于构建通用药物设计平台，通过 AI 模拟代谢过程来加速审批、跳过部分动物实验。但监管制度的进化速度是关键瓶颈——技术突破需要与制度改革同步推进。

%% ==================== Section 6 ====================
